\documentclass[11pt]{article}
\usepackage[T1]{fontenc}
\usepackage{lmodern}
\usepackage[a4paper,margin=2.3cm]{geometry}
\usepackage{mriseq}
\usepackage{booktabs,array,tabularx,xltabular}
\usepackage[most]{tcolorbox}
\usepackage{subcaption}
\usepackage[hidelinks]{hyperref}

\tcbset{
  mriseqex/.style={
    enhanced, breakable, colback=white, colframe=black!20, boxrule=0.4pt, arc=1mm,
    listing options={basicstyle=\ttfamily\small, columns=fullflexible, keepspaces,
      keywordstyle=\color{black}, commentstyle=\color{black!55}},
    left=4pt, right=4pt, segmentation style={black!20}}}
\newtcblisting{example}{mriseqex}
\newtcblisting{exampleonly}{mriseqex,listing only}

\newcommand\pkg[1]{\textsf{#1}}
\newcommand\cs[1]{\texttt{\textbackslash#1}}
\newcommand\key[1]{\texttt{#1}}
\newcommand\meta[1]{\ensuremath{\langle}\textit{#1}\ensuremath{\rangle}}
\newcommand\marg[1]{\texttt{\{}\meta{#1}\texttt{\}}}
\newcommand\oarg[1]{\texttt{[}\meta{#1}\texttt{]}}

\title{\pkg{mriseq}: MRI pulse sequence diagrams with Ti\emph{k}Z\\[2pt]
  \large Version 0.3.0}
\author{Gabriel Lindeblad\\\small\url{https://github.com/Glindeb/mriseq}}
\date{2026-11-17}

\begin{document}
\sloppy
\maketitle

\begin{abstract}\noindent
\pkg{mriseq} draws magnetic resonance imaging (MRI) pulse sequence diagrams as
vector graphics directly inside a \LaTeX{} document. It provides physically
consistent ready-made spin-echo, gradient-echo, inversion-recovery, echo-planar
(EPI) and turbo spin-echo diagrams
computed from imaging parameters (TE, TR, TI, slice thickness, FOV, matrix,
dwell time, ...), and a small set of building blocks (RF pulses, trapezoid
gradients, phase-encoding tables, ADC windows, signals, free waveforms, markers
and timing intervals) to draw any other sequence or any part of one.
What is shown is configurable: the rows, the time window, which timings are
marked, and whether numbers and a time axis are printed. It ships with a colour
style, a black-and-white style and a grayscale style, and you can define your own.
\end{abstract}

\begin{center}
\begin{minipage}{0.9\linewidth}\small
\textbf{Notice on the use of generative AI.} This package, its documentation and
its examples were developed with substantial assistance from a generative AI tool
(Claude Code, by Anthropic), under the direction of the author. The author is
responsible for the content. Please check the diagrams you publish and report
errors at \url{https://github.com/Glindeb/mriseq/issues}.
\end{minipage}
\end{center}

\tableofcontents

% =====================================================================
\section{Quick start}

\begin{exampleonly}
\usepackage{mriseq}            % or \usepackage[bw]{mriseq}
...
\begin{figure}
  \centering
  \SpinEchoDiagram
  \caption{Spin-echo pulse sequence.}
\end{figure}
\end{exampleonly}

\begin{example}
\SpinEchoDiagram
\end{example}

Add \key{detailed} for numbers, a time axis in ms and a light grid. The
numbers are computed from the imaging parameters (here the defaults:
5\,mm slice, 240\,mm FOV, $64\times64$ matrix, 50\,\textmu s dwell time,
2\,ms sinc pulses with time--bandwidth product 4):
\begin{example}
\SpinEchoDiagram[detailed, rows={rf,gz,gy,gx,adc}]
\end{example}

The same in black and white, with a longer echo time and crusher gradients:
\begin{example}
\SpinEchoDiagram[style=bw, TE=12, crushers, values]
\end{example}

The package requires a \LaTeX{} kernel from 2022-06 or newer and Ti\emph{k}Z.
It works with pdf\LaTeX, Xe\LaTeX{} and Lua\LaTeX. Diagrams take the width of
the surrounding text (\cs{linewidth}), so they also work in \texttt{minipage}s
and \texttt{subfigure}s; see Section~\ref{sec:size}.

% =====================================================================
\section{Ready-made sequences}

\begin{center}
\begin{tabular}{@{}ll@{}}
\toprule
Stand-alone command & inside a \texttt{mriseq} environment \\
\midrule
\cs{SpinEchoDiagram}\oarg{keys}          & \cs{spinecho}\oarg{keys} \\
\cs{GradientEchoDiagram}\oarg{keys}      & \cs{gradientecho}\oarg{keys} \\
\cs{InversionRecoveryDiagram}\oarg{keys} & \cs{inversionrecovery}\oarg{keys} \\
\cs{EPIDiagram}\oarg{keys}               & \cs{epi}\oarg{keys} \\
\cs{TurboSpinEchoDiagram}\oarg{keys}     & \cs{turbospinecho}\oarg{keys} \\
\cs{DiffusionSpinEchoDiagram}\oarg{keys} & \cs{diffusionspinecho}\oarg{keys} \\
\cs{DiffusionEPIDiagram}\oarg{keys}      & \cs{diffusionepi}\oarg{keys} \\
\bottomrule
\end{tabular}
\end{center}

The stand-alone commands accept \emph{both} sequence keys
(Section~\ref{sec:seqkeys}) and diagram keys (Section~\ref{sec:diagkeys}) in
one list. The lower-case forms add the elements of the sequence to a
\texttt{mriseq} environment, so you can combine them with your own elements and
annotations (Section~\ref{sec:combine}).

\subsection{Gradient echo}
\begin{example}
\GradientEchoDiagram[flip label=$\alpha$]
\end{example}

With a spoiler, a phase rewinder, the next repetition (the long, empty gap
is compressed automatically, see \key{compress}) and values:
\begin{example}
\GradientEchoDiagram[spoiler, rewinder, show TR, TR=30, values,
                     phase highlight=0.5]
\end{example}

\subsection{Inversion recovery}
The default read-out is a spin echo; \key{readout=gre} gives a gradient-echo
read-out. The $M_z$ row shows the longitudinal magnetisation computed from
\key{T1} (it is not shown by default; add it to \key{rows}).
\begin{example}
\InversionRecoveryDiagram[rows={rf,gz,gy,gx,sig,mz}, TI=300, T1=600]
\end{example}

Note how the signal is inverted: at $\mathrm{TI}=300$\,ms and
$T_1=600$\,ms the longitudinal magnetisation is still negative,
$1-2e^{-\mathrm{TI}/T_1}<0$, when the $90^\circ$ pulse is applied.

\subsection{Echo-planar imaging}
\cs{EPIDiagram} draws a single-shot, blipped EPI read-out: after the excitation the
read and phase prephasers move to a corner of $k$-space, then the read-out
gradient alternates sign for every $k$-space line while small triangular blips on
$G_y$ step to the next line. TE is the time to the centre line of $k$-space.
The whole echo train is drawn to scale, without an axis cut.
EPI uses its own defaults: \key{phase steps=16} (16 lines keep the diagram
readable; use \key{phase steps=64} for a realistic train), \key{dwell=0.005}
(200\,kHz), \key{grad ramp=0.1} and \key{grad scale=row} (each gradient row is
scaled to its own largest amplitude, so the small blips stay visible; use
\key{grad scale=shared} for one common scale).
\begin{example}
\EPIDiagram
\end{example}
\key{refocusing} gives a spin-echo EPI; \key{kspace} adds the $k$-space
trajectory (Section~\ref{sec:kspace}):
\begin{example}
\EPIDiagram[refocusing, annotate={TE/2,TE}, values, kspace,
            rows={rf,gz,gy,gx,adc}]
\end{example}
Zoom into a few lines of the train with a time window:
\begin{example}
\EPIDiagram[detailed, rows={gy,gx,adc}, time=4:8]
\end{example}

\paragraph{Segmented (multi-shot) EPI.} \key{shots=\meta{n}} acquires every
$n$-th line per shot: the blips are $n$ lines high and the phase prephaser of
shot $s$ starts $s$ lines further in. By default two shots are drawn
(\key{repetitions} changes that); the $k$-space panel shows the interleaving:
\begin{example}
\EPIDiagram[shots=4, repetitions=4, TR=60, rows={rf,gy,gx}, kspace]
\end{example}

\subsection{Turbo spin echo}
\cs{TurboSpinEchoDiagram} draws a CPMG echo train of \key{echo train length}
(\key{ETL}, default 4) refocusing pulses with spacing \key{echo spacing}
(\key{ESP}, \texttt{auto} is the shortest possible, rounded up to whole ms).
Every echo has its own phase encoding and rewinder; the highlighted line shows how
$k_y$ is stepped through the train. The echo with $k_y=0$ is
\key{effective echo} (default: the middle one) and gives the effective echo time
$\mathrm{TE}_\mathrm{eff}$.
\begin{example}
\TurboSpinEchoDiagram[rows={rf,gz,gy,gx,sig}]
\end{example}

\subsection{Diffusion weighting}
\cs{DiffusionSpinEchoDiagram} and \cs{DiffusionEPIDiagram} add a
Stejskal--Tanner pair of diffusion gradients around the $180^\circ$ pulse, on
the row \texttt{gd} ($G_\mathrm{D}$). From the b-value \key{b} (s/mm$^2$), the
gradient amplitude \key{diffusion amp} (mT/m) and ramp time \key{diffusion ramp}
($\varepsilon$) the lobe duration $\delta$ is solved from
\[
  b = \gamma^2 G^2\Bigl[\delta^2\bigl(\Delta-\tfrac{\delta}{3}\bigr)
      + \tfrac{\varepsilon^3}{30} - \tfrac{\delta\varepsilon^2}{6}\Bigr]
\]
with the lobes placed right next to the refocusing pulse; TE grows accordingly.
\key{diffusion axis=\{gx,gy,gz\}} puts the lobes on the imaging axes instead
(each with amplitude $G/\sqrt{n}$). The echo is weighted by
$e^{-b D}$ with \key{diffusivity} $D$ (mm$^2$/s). Diffusion diagrams scale each
gradient row on its own (\key{grad scale=row}).
\begin{example}
\DiffusionSpinEchoDiagram[values, rows={rf,gz,gx,gd,sig}]
\end{example}
\begin{example}
\DiffusionEPIDiagram[b=500]
\end{example}

\subsection{Several repetitions}
\key{repetitions=\meta{n}} draws $n$ complete repetitions, \key{TR} apart (default
\texttt{auto}: one, or two for segmented EPI). The
empty time between them is cut (see \key{compress}), every TR is marked, the
highlighted phase-encoding line steps from one repetition to the next, and the
$M_z$ row shows how the magnetisation approaches its steady state.
\key{show TR} adds the start of one more repetition.
\begin{example}
\SpinEchoDiagram[repetitions=3, TR=600, rows={rf,gz,gy,gx,sig,mz}]
\end{example}
\begin{example}
\GradientEchoDiagram[repetitions=3, TR=12, spoiler, ticks,
                     rows={rf,gz,gy,gx,mz}]
\end{example}
\begin{example}
\EPIDiagram[repetitions=2, ticks, break style=zigzag,
            rows={rf,gz,gy,gx,adc}]
\end{example}

\subsection{Sequence keys}\label{sec:seqkeys}
All times are in ms, gradients in mT/m, areas in mT\,ms/m. The keys can also be
set globally with \cs{mriseqset}.

{\small
\begin{xltabular}{\linewidth}{@{}>{\ttfamily\raggedright\arraybackslash}p{0.28\linewidth}>{\ttfamily}lX@{}}
\toprule
\normalfont Key & \normalfont Default & Meaning \\
\midrule
\endhead
TE & auto & echo time; \texttt{auto} is the shortest possible TE rounded up to whole ms. A warning is issued if a given TE is too short. \\
TR & 500 & repetition time (only drawn with \key{show TR}) \\
TI & 400 & inversion time (inversion recovery) \\
show TR & false & draw the start of the next repetition and the TR interval \\
repetitions & auto & number of complete repetitions drawn \\
mark TR & true & mark TR between repetitions \\
refocusing & false & EPI: spin-echo EPI with a $180^\circ$ pulse at TE/2 \\
echo train length, ETL & 4 & TSE: number of echoes \\
echo spacing, ESP & auto & TSE: echo spacing \\
effective echo & auto & TSE: echo with $k_y=0$ (the middle one) \\
shots & 1 & EPI: number of shots (segments) \\
b & 0 / 1000 & b-value in s/mm$^2$ (diffusion diagrams) \\
diffusion amp & 40 & diffusion gradient amplitude (mT/m) \\
diffusion ramp & 0.3 & ramp time of the diffusion gradients \\
diffusion axis & gd & row(s) of the diffusion gradients \\
diffusivity & 0.0008 & apparent diffusion coefficient (mm$^2$/s) for the signal \\
delta label, Delta label & & labels of the diffusion intervals ($\delta$, $\Delta$) \\
equal areas & false & shade the pairs of equal gradient areas (SE, GRE; Section~\ref{sec:shade}) \\
flip & 90 / 30 & excitation flip angle in degrees (SE, IR / GRE) \\
flip label & & text above the excitation pulse instead of the angle, e.g.\ \verb|$\alpha$| \\
excitation phase & & subscript of the excitation angle, e.g.\ \texttt{x} gives $90^\circ_x$ \\
refocusing phase & & same for the $180^\circ$ pulse \\
readout & se & read-out of inversion recovery: \texttt{se} or \texttt{gre} \\
pulse dur & 2 & duration of the excitation pulse \\
refocusing dur & 2 & duration of the refocusing pulse \\
inversion dur & 2 & duration of the inversion pulse \\
pulse tbw & 4 & time--bandwidth product of the sinc pulses \\
slice & 5 & slice thickness in mm \\
fov & 240 & field of view in mm \\
matrix & 64 & sets both \key{read points} and \key{phase steps} \\
read points & 64 & number of read-out samples \\
phase steps & 64 / 16 & number of phase-encoding steps (EPI: 16) \\
dwell & 0.05 / 0.005 & ADC dwell time (ms per sample; EPI: 0.005) \\
grad ramp & 0.2 / 0.1 & ramp time of all gradients (EPI: 0.1) \\
prephaser dur & 1 & duration of the slice rephaser, phase encoder and read prephaser (incl.\ ramps) \\
lead & 0.5 & time before the first gradient starts \\
crushers & false & crusher gradients around the $180^\circ$ pulse (SE) \\
crusher dur & 0.8 & duration of each crusher \\
spoiler & false & spoiler gradient on $G_z$ after the read-out (GRE) \\
inversion spoiler & true & spoiler after the inversion pulse (IR) \\
spoiler dur & 1.4 & spoiler duration \\
rewinder & false & phase rewinder after the read-out (GRE) \\
phase lines & 9 & number of lines drawn in the phase-encoding table \\
phase highlight & & draw one phase step in the accent colour, value in $[-1,1]$ \\
annotate & \normalfont see text & timing intervals to draw: any of \texttt{TE}, \texttt{TE/2}, \texttt{TI}, \texttt{TR}, \texttt{ESP} (EPI, TSE), \texttt{train} (EPI), \texttt{delta}, \texttt{Delta} (diffusion). Defaults: SE \texttt{\{TE/2,TE\}}, GRE and EPI \texttt{\{TE\}}, IR \texttt{\{TI,TE\}}, TSE \texttt{\{ESP,TE\}}, diffusion \texttt{\{TE,delta,Delta\}}. Use \texttt{annotate=\{\}} for none. \\
TE label, TE/2 label, TR label, TI label, ESP label & \normalfont\ttfamily TE, \dots & the text of the intervals \\
pulse markers & true & dashed lines through all rows at the RF pulse centres \\
echo marker & true & marker at the echo centre \\
echo label & SE / GE & label of the echo marker \\
signal & true & emit the FID and echo for the \texttt{sig} row \\
T1, T2, T2* & 900, 90, 40 & relaxation times used for the signal and $M_z$ \\
fid decay & 0.7 & display decay constant of the FID \\
\bottomrule
\end{xltabular}}


\subsection{How the numbers are computed}\label{sec:physics}
With $\bar\gamma = 42.577$\,kHz/mT (\textsuperscript{1}H):
\begin{itemize}
\item slice-select amplitude $G_\mathrm{ss} = \dfrac{\mathrm{TBW}/\tau}{\bar\gamma\,\Delta z}$
      (pulse duration $\tau$, slice thickness $\Delta z$); the RF pulse lies on the flat top;
\item slice rephaser area $-G_\mathrm{ss}\,(\tau/2 + t_\mathrm{ramp}/2)$;
\item read-out amplitude $G_\mathrm{ro} = \dfrac{1/\Delta t}{\bar\gamma\,\mathrm{FOV}}$
      with dwell time $\Delta t$; the ADC window ($N_\mathrm{read}\,\Delta t$) is the flat top,
      centred on the echo;
\item read prephaser area $\pm G_\mathrm{ro}\,(T_\mathrm{acq}/2 + t_\mathrm{ramp}/2)$:
      the same sign as the read-out for a spin echo (the $180^\circ$ pulse inverts the phase),
      the opposite sign for a gradient echo;
\item maximum phase-encoding area $\dfrac{N_\mathrm{phase}/2}{\bar\gamma\,\mathrm{FOV}}$;
\item spin echo: the $180^\circ$ pulse is at TE/2, the echo at TE (both from the centre
      of the excitation pulse); inversion recovery: TI is measured between the
      centres of the inversion and excitation pulses;
\item EPI: echo spacing $\mathrm{ESP}=N_\mathrm{read}\Delta t + 2t_\mathrm{ramp}$, the
      read-out lobes alternate sign, each blip has the area of one $k$-space line,
      $1/(\bar\gamma\,\mathrm{FOV})$, the prephasers move to
      $(-k_{x,\max},-k_{y,\max})$, and TE is the time to line $N_\mathrm{phase}/2$
      (the centre of $k$-space);
\item segmented EPI with $n$ shots: blips of $n$ lines, shot $s$ starts at
      $k_y = -k_{y,\max} + s\,\Delta k_y$;
\item diffusion: Stejskal--Tanner equation for trapezoidal lobes (above), with
      $\gamma = 2\pi\cdot 42.577$\,MHz/T;
\item TSE: refocusing pulses at $(n-\tfrac12)\,\mathrm{ESP}$, echoes at
      $n\,\mathrm{ESP}$, phase encoding before and rewinding after every echo,
      $\mathrm{TE}_\mathrm{eff}$ is the time of the echo with $k_y=0$;
\item the RF amplitude is proportional to flip angle\,/\,(duration $\times$ mean of the
      pulse shape), i.e.\ the $B_1$ amplitude that gives that flip angle;
\item the $M_z$ row follows $M_z(t) = 1-(1-M_z(t_0))\,e^{-(t-t_0)/T_1}$ between pulses,
      and $M_z \to M_z\cos\alpha$ at each pulse;
\item signals are schematic: FID amplitude $\propto\sin\alpha$, echo amplitude
      $\propto e^{-\mathrm{TE}/T_2}$ (SE) or $e^{-\mathrm{TE}/T_2^*}$ (GRE), multiplied by
      $1-2e^{-\mathrm{TI}/T_1}$ for inversion recovery; the shape of the echo is
      illustrative.
\end{itemize}
All rows of one amplitude group share one scale (all gradient rows share the
gradient scale, except with \key{grad scale=row}), so the heights can be compared
between rows.

% =====================================================================
\section{Choosing what is shown}\label{sec:diagkeys}

\subsection{Rows and time window}
\key{rows} selects the rows and their order. Elements on rows that are not
shown are skipped, but still count for the amplitude scale, so a partial diagram
uses the same scale as the full one. \key{time=\meta{from}:\meta{to}} (or
\key{tmin}, \key{tmax}) selects a time window in ms; everything is clipped to it.
Together they let you isolate a part of a sequence, for example the slice
selection of the spin echo, or only the RF and read-out rows:
\begin{example}
\SpinEchoDiagram[rows={rf,gz}, time=0:4.5, detailed, width=0.6\linewidth,
                 annotate={}, markers=false]
\end{example}
\begin{example}
\SpinEchoDiagram[rows={rf,gx,adc}, values, annotate={TE}]
\end{example}

The predefined rows are \texttt{rf}, \texttt{gz}, \texttt{gy}, \texttt{gx},
\texttt{gd} (diffusion), \texttt{adc}, \texttt{sig} (signal), \texttt{mz} and
\texttt{kx}, \texttt{ky}, \texttt{kz} ($k$-space); the default is
\texttt{rows=\{rf,\allowbreak gz,\allowbreak gy,\allowbreak gx,\allowbreak adc,\allowbreak sig\}}. See Section~\ref{sec:rows} for your own rows.

\subsection{Numbers, axis and markers}
{\small
\begin{xltabular}{\linewidth}{@{}>{\ttfamily}l>{\ttfamily}lX@{}}
\toprule
\normalfont Key & \normalfont Default & Meaning \\
\midrule
\endhead
schematic & & preset: no numbers, no axis, no grid (the default state) \\
detailed  & & preset: \key{values}, \key{ticks}, \key{grid} \\
values    & false & sets \key{amp labels}, \key{area labels} and \key{adc labels}; ready-made sequences also print the length of intervals (TE\,=\,7\,ms) \\
amp labels & false & gradient amplitudes with unit (\key{grad unit=mT/m}) \\
area labels & false & gradient areas (\key{area unit} is empty by default) \\
adc labels & false & ADC samples $\times$ dwell time \\
flip labels & true & flip angles above RF pulses \\
precision & 2 & decimals of printed numbers \\
ticks & false & time axis with ticks and numbers (in ms, \key{axis label=ms}) \\
tick step & auto & distance of the ticks \\
axis & false & a plain time arrow labelled \key{time label=\$t\$} (when \key{ticks} is false) \\
grid & false & light vertical lines at the ticks \\
markers & true & show markers (dashed lines at pulses and echoes) \\
row lines & \normalfont style & thin separators between rows \\
baselines & true & zero line of each row \\
row labels & true & row names at the left (column width \key{label width=1.3cm}) \\
\bottomrule
\end{xltabular}}

\paragraph{Where numbers go.} Amplitudes are printed beyond the peak of a lobe
(above a positive, below a negative lobe). Areas are printed inside the flat top
when they fit there completely, otherwise on the other side of the zero line.
Labels on the background get a backing in the background colour
(\key{background color}, white), so that dashed lines are cleanly interrupted.
A number that would overlap another number or another waveform is left out
(a note is written to the log file); labels at the edge of the time window are
moved inwards. Zoom in with a time window to see all numbers of a crowded
sequence, or switch the check off with \key{avoid overlaps=false}.

Numbers are printed with \cs{mriseqformatnumber}, which you can redefine, e.g.\
with \pkg{siunitx}:
\begin{verbatim}
\renewcommand\mriseqformatnumber[1]{\num{#1}}
\end{verbatim}
A key value that contains a comma or an equals sign must be put in braces,
e.g.\ \verb|label={$k=0$}|.

\subsection{Compressing long pauses}
Empty stretches longer than \key{compress threshold} (8\,ms) are compressed
into a break of width \key{break width} (0.5\,cm), drawn as $/\!/$ on each
baseline and on the axis; the time axis continues with the true times after the
break. Switch it off with \key{compress=false}; add breaks yourself with
\cs{gap}. \key{break margin} (0.6\,ms) is the empty time kept on both sides.
\key{break style=slash|zigzag|none} chooses the cut mark.

\subsection{Size, subfigures and narrow diagrams}\label{sec:size}
\key{width} (default \cs{linewidth}) is the total width including the row
labels. Inside a \texttt{subfigure} or \texttt{minipage} \cs{linewidth} is the
width of that box. Diagrams narrower than \key{auto compact} (9\,cm) switch to
the \key{compact} preset automatically: smaller fonts, lower rows and a
narrower label column (set \key{auto compact=0pt} to switch this off, or use
\key{compact} for any diagram). Your own keys are applied on top of it.
\begin{exampleonly}
\begin{figure}
  \begin{subfigure}{0.48\linewidth}
    \SpinEchoDiagram[rows={rf,gz,gy,gx,sig}]
    \caption{Spin echo}
  \end{subfigure}\hfill
  \begin{subfigure}{0.48\linewidth}
    \GradientEchoDiagram[rows={rf,gz,gy,gx,sig}]
    \caption{Gradient echo}
  \end{subfigure}
\end{figure}
\end{exampleonly}
\begin{figure}[h]
  \begin{subfigure}{0.48\linewidth}
    \SpinEchoDiagram[rows={rf,gz,gy,gx,sig}]
    \caption{Spin echo}
  \end{subfigure}\hfill
  \begin{subfigure}{0.48\linewidth}
    \GradientEchoDiagram[rows={rf,gz,gy,gx,sig}]
    \caption{Gradient echo}
  \end{subfigure}
\end{figure}

Further size keys: Alternatively \key{xscale=1cm} fixes the length of one ms.
\key{row scale} multiplies all row heights, \key{amp fill} (0.85) is the
fraction of the half row height used by the largest amplitude, \key{grad max},
\key{rf max}, \key{signal max} fix the amplitude that fills a row instead of
using the largest one. \key{grad scale=row} scales each gradient row to its own
largest amplitude instead of one scale for all gradient rows (\key{shared}, the
default except for EPI). \key{tikz=\{...\}} passes options to the
\texttt{tikzpicture}.

% =====================================================================
\section{\texorpdfstring{$k$}{k}-space}\label{sec:kspace}
The $k$-space position $k(t) = \bar\gamma\int G\,dt$ is computed exactly from
the gradients on the rows \texttt{gx}, \texttt{gy} and \texttt{gz}: an
excitation pulse sets $k$ to zero at its centre, a refocusing or inversion pulse
mirrors it ($k\to-k$). A phase-encoding table contributes its highlighted step
(\key{phase highlight}, \key{highlight}) or else its outermost step.

\paragraph{$k(t)$ rows.} Add \texttt{kx}, \texttt{ky} and/or \texttt{kz} to
\key{rows}:
\begin{example}
\GradientEchoDiagram[rows={rf,gy,gx,kx,ky}, phase highlight=0.5]
\end{example}

\paragraph{Trajectory panel.} \key{kspace=right|left|below} draws the
$k_x$--$k_y$ trajectory next to the diagram, \key{kspace size} (3.6\,cm) wide.
The thin line is the path, the thick line the part sampled by the ADC, the
dotted line the jump at a refocusing pulse; the dot marks $k=0$. Only the part
inside the time window is drawn, so with \key{time} you can show how a part of
the sequence moves through $k$-space. The axes are scaled to the largest
$|k_x|$ and $|k_y|$ (\key{kspace isotropic} uses one scale for both);
\key{kspace title} adds a title.
\begin{example}
\SpinEchoDiagram[rows={rf,gz,gy,gx}, kspace=left, phase highlight=0.75]
\end{example}

% =====================================================================
\section{Highlighting areas}\label{sec:shade}
\cs{shadearea}\oarg{keys}\marg{row}\marg{t1}\marg{t2} fills the part of the
gradients of one row between two times. \key{style=A|B|C|D} picks one of four
matching fills (colours in the colour style, hatchings in the black-and-white
style) or any Ti\emph{k}Z fill options; \key{label} adds a label inside the
area, \key{label=auto} prints its area. The ready-made SE and GRE diagrams mark
the classic pairs of equal areas with \key{equal areas}: the second half of the
slice-select gradient and the slice rephaser (A), and the read prephaser and the
first half of the read-out (B).
\begin{example}
\GradientEchoDiagram[rows={rf,gz,gx}, equal areas]
\end{example}
\begin{example}
\begin{mriseq}[rows={rf,gx}, ticks, style=bw]
  \rfpulse{center=1, dur=1, flip=90}
  \gradient{gx}{start=2, flat=0.6, amp=-1}
  \gradient{gx}{start=3, flat=3, amp=0.5}
  \shadearea[label=auto]{gx}{2}{3}
  \shadearea[label=auto]{gx}{3}{4.6}
  \shadearea[style=B]{gx}{4.6}{6.4}
\end{mriseq}
\end{example}

% =====================================================================
\section{Drawing your own sequences}\label{sec:own}

\begin{exampleonly}
\begin{mriseq}[<diagram keys>]
  <elements>
\end{mriseq}
\end{exampleonly}

Elements are collected first and drawn at the end of the environment, so the
order in which you write them does not matter (except for overlaps). Every
element takes a key list. Times are in ms; amplitudes are in any unit you like
(each amplitude group is scaled to its largest value).

\begin{example}
\begin{mriseq}[rows={rf,gz,gx,sig}, ticks, grad labels=plain]
  \rfpulse{center=1.5, dur=1, flip=90, phase=x}
  \gradient{gz}{center=1.5, flat=1, amp=1}
  \gradient{gz}{start=2.2, flat=0.4, area=-0.6}
  \gradient{gx}{start=2.2, flat=0.6, amp=-1, label=prephaser}
  \gradient{gx}{center=5, flat=2.4, amp=0.8}
  \fid{start=2, t2=0.4}
  \echo{center=5, t2=0.4, cycles=2}
  \marker[kind=echo, label=echo]{5}
  \interval{1.5}{5}{TE}
\end{mriseq}
\end{example}

\subsection{Element commands}
\begin{center}\small
\begin{tabularx}{\linewidth}{@{}lX@{}}
\toprule
Command & Element \\
\midrule
\cs{rfpulse}\oarg{row}\marg{keys}         & RF pulse (row \texttt{rf} by default) \\
\cs{gradient}\marg{row}\marg{keys}        & trapezoid gradient \\
\cs{phasetable}\oarg{row}\marg{keys}      & phase-encoding table (row \texttt{gy}) \\
\cs{readout}\oarg{row}\marg{keys}         & ADC window (row \texttt{adc}) \\
\cs{fid}\oarg{row}\marg{keys}             & free induction decay (row \texttt{sig}) \\
\cs{echo}\oarg{row}\marg{keys}            & echo signal (row \texttt{sig}) \\
\cs{waveform}\marg{row}\oarg{keys}\marg{t/a, t/a, ...} & piecewise-linear waveform \\
\cs{marker}\oarg{keys}\marg{t}            & vertical line through all rows \\
\cs{interval}\oarg{keys}\marg{t1}\marg{t2}\marg{label} & timing interval with arrows \\
\cs{gap}\marg{t1}\marg{t2}                & compress $[t_1,t_2]$ into a break \\
\cs{shadearea}\oarg{keys}\marg{row}\marg{t1}\marg{t2} & shade gradient areas, see \ref{sec:shade} \\
\cs{annotate}\marg{TikZ code}             & arbitrary Ti\emph{k}Z code, see \ref{sec:cs} \\
\bottomrule
\end{tabularx}
\end{center}

\paragraph{Timing keys (all waveform elements).}
\key{start}, \key{center} (alias \key{at}) or \key{end} place the element;
for gradients \key{center} is the centre of the flat top. \key{dur} is the
total duration.

\paragraph{\cs{rfpulse}.}
\key{flip} (degrees; also the default label and amplitude), \key{amp}
(overrides the amplitude), \key{dur} (default \key{rf dur=2}),\\
\key{shape=sinc|gauss|hard|sech} (default \key{rf shape}), \key{tbw}
(sinc time--bandwidth product, default 4),\\ \key{window=hamming|hann|none}
(default \key{rf window=hamming}),\\ \key{phase} (subscript), \key{label},
\key{role=excite|refocus|invert} (colour; default: refocus for flip $>135^\circ$,
otherwise excite).

\paragraph{\cs{gradient} and \cs{phasetable}.}
\key{amp}, \key{flat}, \key{ramp} (both ramps, default \key{ramp=0.2}),
\key{ramp up}, \key{ramp down}, \key{area}. Give any two of \key{amp},
\key{flat}, \key{area} (or \key{dur} instead of \key{flat}). \key{label}
replaces the amplitude label; \key{split=\meta{t}} prints the area before and
after time $t$ (e.g.\ the centre of an RF pulse). Phase tables also take
\key{lines} (default 9), \key{arrow=down|up|none}, \key{highlight} (a value in
$[-1,1]$) and \key{fan} (draw each step as a full trapezoid).

\paragraph{\cs{readout}.} \key{samples}, \key{dwell} (gives the duration if
\key{dur} is not set), \key{kzero=\meta{t}} (dot at the centre of $k$-space),
\key{sample ticks} (true), \key{label}.

\paragraph{\cs{fid} and \cs{echo}.} \key{amp}, \key{t2} (decay constant),
\key{cycles} (oscillations per decay constant), \key{oscillate} (true; false
draws the envelope), \key{resolution}.

\paragraph{\cs{waveform}.} \key{filled} (true), \key{label}, and
\key{group} to put it in another amplitude group.

\paragraph{Common keys.} \key{color} (line colour), \key{fill}, \key{style}
(extra Ti\emph{k}Z options), \key{values}, \key{amp labels}, \key{area labels}
to override the diagram setting for one element, \key{occupy=false} to exclude
an element from the time range and from the automatic compression.

\subsection{Markers and intervals}
\cs{marker} takes \key{kind=plain|echo|excite|refocus|invert} (colour and look),
\key{label} and \key{color}. \cs{interval} takes \key{pos=top|bottom|\meta{row}},
\key{level} (stacking level above the diagram, 1 is the lowest),
\key{height} (position inside a row, 0 to 1), \key{guides} (dotted lines from
a top interval down to the diagram) and \key{style}.

\begin{example}
\begin{mriseq}[rows={rf,gx,adc}, ticks, compress threshold=5]
  \rfpulse{center=1, dur=1, flip=90}
  \gradient{gx}{start=2, flat=1, amp=-1}
  \gradient{gx}{center=6, flat=3, amp=0.67}
  \readout{center=6, samples=32, dwell=0.09375, kzero=6}
  \interval[level=1]{1}{6}{$T_E$}
  \interval[pos=gx, height=0.85, guides=false]{4.5}{7.5}{$T_\mathrm{acq}$}
  \interval[pos=bottom]{0.5}{1.5}{RF}
  \marker[kind=echo, label={$k=0$}]{6}
  \gap{8}{40}
  \rfpulse{center=41, dur=1, flip=90}
  \interval[level=2]{1}{41}{$T_R$}
\end{mriseq}
\end{example}

\subsection{Combining ready-made sequences and own elements}\label{sec:combine}
\begin{example}
\begin{mriseq}[rows={rf,gz,gy,gx,sig}]
  \spinecho[annotate={TE}, echo marker=false]
  \marker[kind=echo, label={Hahn echo}]{8.7}
  \annotate{\node[font=\scriptsize, anchor=west]
     at (seq cs:t=9.4, row=sig, a=0.6) {$\propto e^{-\mathrm{TE}/T_2}$};}
\end{mriseq}
\end{example}

\subsection{Ti\emph{k}Z coordinates}\label{sec:cs}
Inside \cs{annotate} the coordinate system \texttt{seq} converts sequence
coordinates to the picture:
\begin{center}
\verb|(seq cs: t=|\meta{ms}\verb|, row=|\meta{row}\verb|, a=|\meta{amplitude}\verb|)|
\end{center}
\key{a} is in the units of that row (after scaling). Without \key{row},
\key{y=\meta{length}} gives an absolute height (0 is the top of the first row).

\begin{example}
\begin{mriseq}[rows={rf,gx}]
  \rfpulse{center=1, dur=1.2, flip=90}
  \gradient{gx}{start=2, flat=0.6, amp=-1}
  \gradient{gx}{start=3, flat=2.4, amp=0.5}
  \annotate{
    \fill[red!25, opacity=0.6] (seq cs:t=2.2,row=gx,a=0)
       rectangle (seq cs:t=2.8,row=gx,a=-1);
    \draw[red, <-] (seq cs:t=2.5,row=gx,a=-0.5) -- ++(1,-0.2)
       node[right, font=\footnotesize]{area $=-A$};}
\end{mriseq}
\end{example}

% =====================================================================
\section{Rows}\label{sec:rows}
\cs{mriseqdefinerow}\marg{name}\marg{keys} defines or changes a row. Keys:
\key{label}, \key{sublabel}, \key{height} (a length), \key{baseline}
(position of the zero line as a fraction of the height; 0.5 for gradients,
0.22 for RF) and \key{group} (amplitude group: \texttt{rf}, \texttt{grad},
\texttt{adc}, \texttt{sig}, \texttt{mz} or any new name). Inside a diagram's
options use \key{row=\{\meta{name}\}\{\meta{keys}\}}.

\key{grad labels=physical|logical|plain|words} switches the gradient labels
between $G_z$/$G_y$/$G_x$ with sublabels, $G_\mathrm{SS}$/$G_\mathrm{PE}$/$G_\mathrm{RO}$,
plain $G_z$/$G_y$/$G_x$, and $G_\mathrm{slice}$/$G_\mathrm{phase}$/$G_\mathrm{read}$.

\begin{example}
\mriseqdefinerow{gs}{label=$G_\mathrm{s}$, sublabel=spoiler, group=grad}
\begin{mriseq}[rows={rf,gs,gx}, grad labels=logical, row={rf}{label=$B_1$}]
  \gradientecho[signal=false]
  \gradient{gs}{start=8, flat=1.5, amp=12}
\end{mriseq}
\end{example}

% =====================================================================
\section{Styles}
\subsection{Built-in styles}
\key{style=color} (default), \key{style=bw} (all lines and text black, filled
areas gray) and \key{style=grayscale} (gray tones for the different pulse
types). Choose one for the document with the package option
(\verb|\usepackage[bw]{mriseq}|) or \verb|\mriseqset{style=bw}|, or for one
diagram with its \key{style} key.

\begin{example}
\GradientEchoDiagram[style=bw, rows={rf,gz,gy,gx,sig}, spoiler]
\end{example}
\begin{example}
\InversionRecoveryDiagram[style=grayscale, rows={rf,gz,gy,gx,adc}]
\end{example}

\subsection{Appearance keys}
\begin{center}\small
\begin{tabularx}{\linewidth}{@{}>{\ttfamily\raggedright\arraybackslash}p{0.45\linewidth}X@{}}
\toprule
\normalfont Key & Meaning \\
\midrule
rf excite color, rf refocus color, rf invert color, rf color & RF pulses by role (\key{rf color} sets all three) \\
grad color, phase color, adc color, signal color & lines of the elements; \key{phase color} is the accent for phase tables \\
echo color, marker color, interval color & markers and intervals \\
text color, sublabel color, value color, axis color & text \\
baseline color, grid color, row line color, background color & background (\key{background color} is the backing of labels) \\
k color & $k$-space trajectory \\
area A style, \dots, area D style & Ti\emph{k}Z fill options of the shaded areas \\
rf fill mix, grad fill mix, adc fill mix, signal fill mix & fill = line colour mixed with this percentage and white \\
rf fill, grad fill, adc fill & a fixed fill colour instead of the mix \\
line width, thin line width & 0.7pt / 0.4pt in the colour style \\
label font, sublabel font, value font, axis font, interval font, flip font & fonts \\
\bottomrule
\end{tabularx}
\end{center}

All drawing is done with Ti\emph{k}Z styles under \texttt{mriseq/} that you can
change with \verb|\tikzset|, e.g.\ \texttt{mriseq/rf}, \texttt{mriseq/grad},
\texttt{mriseq/phase}, \texttt{mriseq/phase lines}, \texttt{mriseq/adc},
\texttt{mriseq/signal}, \texttt{mriseq/marker}, \texttt{mriseq/echo marker},
\texttt{mriseq/interval}, \texttt{mriseq/interval label},
\texttt{mriseq/baseline}, \texttt{mriseq/grid}, \texttt{mriseq/axis},
\texttt{mriseq/row label}, \texttt{mriseq/value}, \texttt{mriseq/value outside},
\texttt{mriseq/flip label}, \texttt{mriseq/area A} (to \texttt{D}),
\texttt{mriseq/kspace path}, \texttt{mriseq/kspace adc}, \texttt{mriseq/kspace frame}.

\subsection{Your own style}
\begin{example}
\mriseqdefinestyle{report}{
  style=color, rf excite color=teal, rf refocus color=orange!80!black,
  grad color=black!70, grad fill mix=10, line width=0.9pt,
  label font=\small\sffamily, value font=\scriptsize\sffamily}
\GradientEchoDiagram[style=report, rows={rf,gz,gy,gx}, values]
\end{example}

% =====================================================================
\section{Compile time}
Every diagram is computed by \TeX{} itself, which takes roughly 0.5--3\,s per
diagram (EPI with a $k$-space panel and long echo trains are the slowest).
Rows that are not shown are not computed where possible (signal, $M_z$, $k$).
For a report with many diagrams, Ti\emph{k}Z's externalisation caches every
diagram as a PDF and only recompiles changed ones; a cached diagram costs
almost nothing (all computation happens inside the \texttt{tikzpicture}, which
is skipped):
\begin{exampleonly}
\usetikzlibrary{external}
\tikzexternalize[prefix=tikzcache/]   % needs -shell-escape
\end{exampleonly}

% =====================================================================
\section{Limitations and plans}
\begin{itemize}
\item Ready-made sequences: SE, GRE, IR (SE or GRE read-out), EPI (GRE or SE,
      single or multi-shot), TSE, diffusion-weighted SE and SE-EPI. Further
      candidates: stimulated echo, balanced SSFP, spiral and radial read-outs.
\item Signal shapes are illustrative (see Section~\ref{sec:physics}).
\item An element that crosses a compressed gap is drawn compressed as well.
\item Overlapping gradients on the same row are drawn on top of each other (for
      $k$-space they are added).
\item Nested \texttt{mriseq} environments are not supported.
\end{itemize}

\section{Change log}
\begin{description}
\item[0.3.0 (2026-11-17)] First public release: SE, GRE, IR, EPI
  (single/multi-shot, GRE/SE), TSE, diffusion SE and SE-EPI; repetitions and
  axis cut marks; $k(t)$ rows and $k$-space panel; area shading; label
  placement without overlaps; compact mode for narrow diagrams.
\end{description}

\end{document}
